\documentclass[12pt]{article}
\usepackage[T1]{fontenc}
\usepackage[utf8]{inputenc}
\usepackage{lmodern}
\usepackage{textcomp}
\usepackage{enumitem}
\usepackage{geometry}
\usepackage{graphicx}
\usepackage{amsmath, amssymb}
\geometry{margin=1in}
\begin{document}
\begin{center}
Chemistry - Graduate Level Assessment 5 \
Total Marks: 40 \
Answer all questions.
\end{center}

\section*{Multiple Choice (10 marks)}
\begin{enumerate}[label=\arabic*.]
\item How many gallons of a liquid that is 74 percent alcohol must be combined with 5 gallons of one that is 90 percent alcohol in order to obtain a mixture that is 84 percent alcohol?
\begin{enumerate}[label=(\alph*)]
    \item 3 gallons
    \item 5 gallons
    \item 8 gallons
    \item 2 gallons
    \item 1 gallon
\end{enumerate}
\item Determine themolarityof a 40.0\% solution ofHClwhich has a density of 1.20 g/ml.
\begin{enumerate}[label=(\alph*)]
    \item 14.15 M
    \item 16.20 M
    \item 15.00 M
    \item 12.75 M
    \item 9.95 M
\end{enumerate}
\item What is the ionization constant, Ka, for a weak monoprotic acid if a 0.060 molar solution has a pH of 2.0?
\begin{enumerate}[label=(\alph*)]
    \item 3.0 \ensuremath{\times} 10^-3
    \item 2.0 \ensuremath{\times} 10^-3
    \item 1.7 \ensuremath{\times} 10^-1
    \item 1.5 \ensuremath{\times} 10^-2
    \item 2.0 \ensuremath{\times} 10^-1
\end{enumerate}
\item A gas at $250 \mathrm{~K}$ and $15 \mathrm{~atm}$ has a molar volume 12 per cent smaller than that calculated from the perfect gas law. Calculate the molar volume of the gas.
\begin{enumerate}[label=(\alph*)]
    \item 1.2$\mathrm{dm}^3 \mathrm{~mol}^{-1}$
    \item 1.35$\mathrm{dm}^3 \mathrm{~mol}^{-1}$
    \item 2.0$\mathrm{dm}^3 \mathrm{~mol}^{-1}$
    \item 0.8$\mathrm{dm}^3 \mathrm{~mol}^{-1}$
    \item 1.8$\mathrm{dm}^3 \mathrm{~mol}^{-1}$
\end{enumerate}
\item 2.3 g of ethanol (C\_2 H\_5 OH, molecular weight = 46 g/mole) is added to 500 g of water. Determine the molality of the resulting solution.
\begin{enumerate}[label=(\alph*)]
    \item 0.12 molal
    \item 0.05 molal
    \item 0.4 molal
    \item 0.2 molal
    \item 0.1 molal
\end{enumerate}
\end{enumerate}

\section*{True / False (10 marks)}
\textit{Answer True or False}
\begin{enumerate}[label=\arabic*.]
\setcounter{enumi}{5}
\item True or False: The root mean square speed of oxygen (O\_2) at the same temperature as hydrogen (H\_2) is 400 m/sec.
\item True or False: The reaction of sulfuric acid with 120 g of metallic calcium produces 33.6 liters of hydrogen gas at standard temperature and pressure (STP).
\item True or False: The pKa of a weak monoprotic acid with a pH of 4.26 in a 0.0025 mol solution is 6.52.
\item True or False: Alkanes are hydrocarbons that do not contain any oxygen atoms in their molecular structure.
\item True or False: In terms of reducing ability, the order is H\_2 \textgreater{} Ni \textgreater{} Co \textgreater{} Al \textgreater{} Na \textgreater{} Ag.
\end{enumerate}

\section*{Long Form Response (20 marks)}
\textit{Answer each question with a clear and complete explanation.}
\begin{enumerate}[label=\arabic*.]
\setcounter{enumi}{10}
\item Given a household cleaning solution with a hydronium ion concentration of 10^-11 M, calculate the pH of the solution and discuss whether it is acidic, basic, or neutral, providing a rationale for your answer.
\item Calculate the cell potential for the reaction Ni + $Sn^2$+ → $Ni^2$+ + Sn, given the concentrations of $Ni^2$+ at 1.3 M and $Sn^2$+ at 1.0 M. Discuss the implications for spontaneity.
\end{enumerate}

\vfill
\noindent\textit{End of Paper}
\end{document}